% ============================================================
% Section 3: Motivating Observations (Medium length for TNNLS)
% Target: ~1.5 pages
% ============================================================

\section{Motivating Observations}
\label{sec:observation}

Before presenting the scoped calculation, we summarize empirical observations that motivate separate analyses of aggregation, feature-only performance, and longer-range label structure. These examples are descriptive and do not establish a universal decision rule.

\subsection{The U-Shape Pattern}

We conduct controlled synthetic experiments on the Contextual Stochastic Block Model (CSBM) with $n=1000$ nodes, varying homophily $h \in \{0.1, 0.2, ..., 0.9\}$ while keeping features constant (MLP accuracy $\approx 99\%$).

\begin{figure}[!t]
\centering
\includegraphics[width=0.95\columnwidth]{u_shape_pub}
\caption{Recorded accuracy across the synthetic homophily sweep. In this configuration, the largest MLP--GNN accuracy gap occurs near $h=0.5$.}
\label{fig:ushape}
\end{figure}

\begin{table}[t]
\centering
\caption{U-Shape Severity by Model Architecture}
\label{tab:ushape}
\begin{tabular}{@{}lcccc@{}}
\toprule
\textbf{Model} & \textbf{Type} & \textbf{Valley} & \textbf{Peak} & \textbf{Amplitude} \\
\midrule
GCN & Replace & $-18.75\%$ & $+0.60\%$ & \textbf{19.35\%} \\
GAT & Replace+Attn & $-5.30\%$ & $+0.55\%$ & 5.85\% \\
GraphSAGE & Concat & $-0.10\%$ & $+0.65\%$ & \textbf{0.75\%} \\
\bottomrule
\end{tabular}
\end{table}

\textbf{Observation 1} (Figure~\ref{fig:ushape}, Table~\ref{tab:ushape}): In this controlled configuration, the recorded GNN curves have a U-shape with a minimum near $h=0.5$. The measured amplitude differs by architecture:
\begin{itemize}
    \item \textbf{GCN}: 19.35 percentage points;
    \item \textbf{GAT}: 5.85 percentage points; and
    \item \textbf{GraphSAGE}: 0.75 percentage points.
\end{itemize}

The ordering is consistent with the hypothesis that retaining an ego channel can reduce sensitivity to harmful neighborhood mixing. The sweep does not isolate concatenation from attention, normalization, parameterization, or optimization, so it cannot establish that design choice as the unique cause.

\subsection{Same Homophily, Opposite Outcomes}

\begin{table}[t]
\centering
\caption{Same Homophily, Different GNN Outcomes}
\label{tab:same_h}
\begin{tabular}{@{}lccccc@{}}
\toprule
\textbf{Dataset} & \textbf{$h$} & \textbf{$k$} & \textbf{MLP} & \textbf{Headroom} & \textbf{GCN Adv.} \\
\midrule
Cora & 0.81 & 3.9 & 75.1\% & 24.9\% & \textbf{+12.7\%} \\
Coauthor-CS & 0.81 & 8.9 & 94.4\% & 5.6\% & $-0.7\%$ \\
\midrule
CiteSeer & 0.74 & 2.7 & 73.8\% & 26.2\% & +3.2\% \\
PubMed & 0.80 & 4.5 & 87.6\% & 12.4\% & +0.2\% \\
\bottomrule
\end{tabular}
\end{table}

\textbf{Observation 2} (Table~\ref{tab:same_h}): Within their recorded protocols, datasets with similar homophily show different GNN outcomes:
\begin{itemize}
    \item Cora ($h=0.81$): GCN gains +12.7\%
    \item Coauthor-CS ($h=0.81$): GCN \textit{loses} 0.7\%
\end{itemize}

The arithmetic positive-gain headroom differs: Cora's MLP accuracy leaves 24.9 percentage points below perfect accuracy, whereas Coauthor-CS leaves 5.6 points. This identity limits positive improvement but does not explain the sign or magnitude of the graph-model gap.

\textbf{Descriptive Observation}: The Cora/Coauthor-CS comparison motivates feature headroom as a diagnostic, but the audited nine-dataset external set does not show a reliable linear correlation between headroom and best-GNN advantage ($r=-0.255$, bootstrap CI crosses zero). We therefore treat this pair as motivation rather than population-level statistical evidence.

\subsection{The Heterophily Paradox}

\textbf{Observation 3}: Heterophilic datasets behave inconsistently in the audited paired protocol:
\begin{itemize}
    \item \textbf{Texas}: H2GCN--MLP is +6.3 percentage points (95\% paired bootstrap interval: 1.3--11.1 points);
    \item \textbf{Wisconsin}: H2GCN--MLP is $-0.8$ points ($-4.3$--2.5), so the interval crosses zero; and
    \item \textbf{Chameleon}: H2GCN--MLP is +14.1 points (12.2--15.9).
\end{itemize}

The same paired protocol does not align these model differences with the ordering of archived two-hop recovery statistics: Wisconsin has positive $\Delta h$ without a clear H2GCN advantage, whereas Chameleon has near-zero $\Delta h$ and a positive H2GCN advantage. Section~\ref{sec:heterophily} therefore treats two-hop recovery as a descriptive mechanism, not a selector.

\subsection{Summary: Three Questions Motivating the Analysis}

These observations reveal three puzzles that homophily alone cannot explain:

\begin{enumerate}
    \item \textbf{U-Shape Puzzle}: Why does GNN performance form a U-shape with homophily? Why is the amplitude architecture-dependent?

    \item \textbf{Same-$h$ Puzzle}: Why do datasets with identical homophily show opposite GNN outcomes?

    \item \textbf{Heterophily Puzzle}: Why does H2GCN clearly exceed MLP on some low-$h$ datasets while matched uncertainty crosses zero on others?
\end{enumerate}

Sections~\ref{sec:snr}--\ref{sec:heterophily} examine three narrower mechanisms without claiming complete answers:
\begin{itemize}
    \item \textbf{U-shape}: the exact fixed-degree ratio is zero at $h=0.5$ ($\rho=0$), giving a mechanism for loss of neighbor-label mean separation under the stated model;
    \item \textbf{Same-$h$}: feature-only accuracy supplies an arithmetic reference but not a structural-utility predictor; and
    \item \textbf{Heterophily}: two-hop recovery describes longer-range label structure but does not determine whether MLP or a particular graph model wins.
\end{itemize}
